\section{Flavoured S-duality}
\label{sec:rains}\label{sec:flavoured}
We now restore the flavour fugacity. The Jacobi elliptic-function
factorisation used in the unflavoured case does not extend directly to
generic $u$. Instead, we recognise the two Schur integrals as
specialisations of a family studied by Rains~\cite{Rains:2018kernel}.
These integrals are taken with respect to the elliptic Selberg measure
\begin{equation}
  \begin{aligned}
    \dd\mu_{p,Q,t}^{(n)}(\mathbf{z})
    ={}&\frac{\bigl((p;p)_\infty(Q;Q)_\infty\Gamma_{p,Q}(t)\bigr)^n}
    {2^n n!}
    \prod_{i<j}
    \frac{\Gamma_{p,Q}(t z_i^{\pm1}z_j^{\pm1})}
    {\Gamma_{p,Q}(z_i^{\pm1}z_j^{\pm1})}\prod_{i=1}^n
    \frac{1}{\Gamma_{p,Q}(z_i^{\pm2})}\,
    \frac{\dd z_i}{2\pi\ii z_i},
  \end{aligned}
  \label{rains:measure}
\end{equation}
where $\Gamma_{p,Q}$ are the elliptic gamma functions and $p$ is an independent
elliptic nome (see Appendix~\ref{app:theta-gamma} for definitions and
identities).

For our purposes the key result is the following.

\begin{theorem}[Rains~\cite{Rains:2018kernel}, Theorem~3.22]
  \label{thm:rains}
  For generic nonzero parameters, the integral
  \begin{equation}
    \mathcal T_n(\mathbf{v};\kappa)=\int
    \mathcal K_\kappa^{(n)}(\mathbf{z};\mathbf{z};t;p,Q)
    \prod_{i=1}^n\prod_{r=1}^4
    \Gamma_{p,Q}\!\left(\sqrt{\frac{pQ}{\kappa}}\,
    v_r^{\pm1}z_i^{\pm1}\right)
    \,\dd\mu_{p,Q,t}^{(n)}(\mathbf{z})
    \label{rains:trace}
  \end{equation}
  satisfies
  \begin{equation}
    \mathcal T_n(\mathbf{v};\kappa)=\mathcal T_n(\mathbf{v}';\kappa),\qquad
    v'_r=\frac{v_r}{\sqrt{v_1v_2v_3v_4}}.
    \label{rains:general-reflection}
  \end{equation}
\end{theorem}

Here $\mathcal K_\kappa^{(n)}(\mathbf{z};\mathbf{x};t;p,Q)$ is an object called
the analytic interpolation kernel, a meromorphic function of
$\mathbf{z}$ and $\mathbf{x}$~\cite[Theorem~3.1]{Rains:2018kernel}.
For a special value of $\kappa$ this kernel can be expressed in terms of
elliptic gamma functions\footnote{\label{rains:contour}This can be found by
  combining Proposition~2.10 of~\cite{Rains:2018kernel} with the corresponding
  analytic continuation result in Theorem~3.1. Integration over $|z_i|=1$
  requires $\left|\sqrt{pQ/\kappa}v_r^{\pm1}\right|<1$, separating poles
  tending zero and infinity.},
\begin{equation}
  \mathcal K_\kappa^{(n)}(\mathbf{z};\mathbf{x};t;p,Q)
  =\prod_{i,j=1}^n\Gamma_{p,Q}(\kappa z_i^{\pm1}x_j^{\pm1}),
  \qquad \kappa=\sqrt{pQ/t}.
  \label{rains:kernel-product}
\end{equation}
Take
\begin{equation}
  p=\frac Qt,\qquad
  v_r=e^{2\ii b_r},\qquad r=1,\ldots,4.
  \label{rains:dictionary}
\end{equation}
Then it is straightforward to show the following via the properties of the
elliptic gamma functions,
\begin{equation}
  \mathcal T_n(v;p)
  =\frac{C_0(Q,t)^n t^{-n^2}}{2^n n!}\,
  \mathcal J_n(\mathbf{b}),
  \label{rains:normalization}
\end{equation}
where
\begin{align}
  \mathcal J_n(\mathbf{b})
  &=\int_{[0,\pi]^n}\prod_{i=1}^n
  W_{\mathbf{b}}(a_i)\frac{\dd a_i}{\pi}\,
  \prod_{i<j}R(a_i-a_j)R(a_i+a_j),
  \label{rains:theta-integral}\\
  W_{\mathbf{b}}(a)
  &=\frac{t}{4q}\,
  \frac{\vartheta_1'(0)^2\vartheta_1(2a)^2}
  {\displaystyle\prod_{r=1}^4
    \vartheta_4(a+b_r)\vartheta_4(a-b_r)}.
  \label{rains:weight}
\end{align}
In this limit Theorem~\ref{thm:rains} becomes
\begin{equation}
  \mathcal J_n(\mathbf{b})=\mathcal J_n(H\mathbf{b}),
  \qquad (H\mathbf{b})_r=b_r-\frac12\sum_{k=1}^4b_k.
  \label{rains:theta-identity}
\end{equation}

This now looks remarkably similar to \eqref{eq:root-integrals}. The only
distinction is the leading weight factor \eqref{rains:weight}. Now consider
the boundary values\footnote{One must be careful in whether the limit is
justified. This is discussed in appendix~\ref{app:rains-limit}.}
\begin{equation}
  \begin{aligned}
    \mathbf{b}_B
    &=\left(c-\frac{\pi\tau}{2},\
    \frac{\pi}{2}-\frac{\pi\tau}{2},\ 0,\ -\frac{\pi}{2}\right).
  \end{aligned}
  \label{rains:additive-BC}
\end{equation}
For these values the weights reduce to $W_{\mathbf{b}_B}(a)=R(a)$. Finally, if
we compute the limiting weight on the reflected side, we find
$W_{H\mathbf{b}_B}(a)=R(2a)$. Thus, Theorem~\ref{thm:rains} reduces to our
sought-after flavoured $S$-duality result.
